The previous assignments focused on uncertainty coming from the model parameters. This assignment shifts to a different source: errors in the input data. Precipitation is the main driver of HBV001a, but it is never measured perfectly. Point gauges pick up undercatch from wind, instrument drift, and siting errors, and converting point measurements to a catchment average adds further uncertainty through the spatial interpolation scheme used. In practice, areal precipitation estimates can carry random errors of several percent at each time step.

How these errors propagate through a nonlinear model is not obvious. A perturbed precipitation series disturbs the water balance at every model storage, and the effect on simulated discharge depends on initial conditions, the current model state, and the parameter values. Of particular interest is what happens to the calibrated parameter set when the model is recalibrated against perturbed inputs: the optimizer will try to compensate for the noise by adjusting parameters, which can push them away from physically plausible values. This raises two specific questions: how much does random precipitation noise degrade model performance when the reference parameters are kept fixed, and how much of that loss can recalibration recover?

\textcite{Oudin.2006} studied these questions for several rainfall-runoff models, including HBV, using both random and systematically biased inputs. They found that random errors are harder to compensate for than bias, and that recalibration restores performance only partially while increasing parameter uncertainty --- especially for parameters controlling slow reservoir dynamics. This assignment replicates the random-perturbation experiment for HBV001a on the study catchment.
